\documentclass[11pt,aspectratio=169]{beamer}

\usetheme{Madrid}
\usecolortheme{seahorse}

\usepackage[utf8]{inputenc}
\usepackage{amsmath,amssymb}
\usepackage{algpseudocode}
\usepackage{algorithmicx}
\usepackage{xcolor}
\usepackage{listings}
\usepackage{booktabs}
\usepackage{tikz}
\usetikzlibrary{arrows.meta,positioning,fit,backgrounds}

\setbeamertemplate{navigation symbols}{}
\setbeamertemplate{footline}[frame number]

\lstdefinestyle{py}{
  language=Python,
  basicstyle=\ttfamily\scriptsize,
  keywordstyle=\color{blue!70!black}\bfseries,
  commentstyle=\color{gray},
  stringstyle=\color{green!40!black},
  showstringspaces=false,
  breaklines=true,
  numbers=left,
  numberstyle=\tiny\color{gray},
  frame=single,
  xleftmargin=1.5em,
  aboveskip=0.3em,
  belowskip=0.3em
}

\title{LeetCode Deep Dive}
\subtitle{From First Instinct to Optimal Algorithm: Greedy Exchange Arguments in Practice\\
\small Week 4 Supplement -- TOAE209/TOAH209: Design and Analysis of Algorithms}
\author{Foreign Trade University}
\date{}

\begin{document}

\begin{frame}
  \titlepage
\end{frame}

\begin{frame}{Outline}
  \tableofcontents
\end{frame}

% =================================================================
\section{A Framework for Attacking a New Greedy Problem}
% =================================================================

\begin{frame}{Why walk through LeetCode problems?}
  \begin{itemize}
    \item This week's \texttt{LEETCODE.md} list gives ten Week 4 problems to
    \emph{practice on your own}. This deck \textbf{narrates the thinking process} for
    two of them, so you can reuse the process on the rest.
    \item We pick one \textcolor{blue!50!black}{\textbf{Medium}} (\emph{Non-overlapping
    Intervals}) that is almost literally this week's lecture, and one
    \textcolor{red!60!black}{\textbf{Hard}} (\emph{Course Schedule III}) that extends
    this week's ``job sequencing with deadlines'' case study with a heap.
  \end{itemize}
\end{frame}

\begin{frame}{The four-step process we will repeat twice}
  \begin{enumerate}
    \item \textbf{First instinct.} Write down the algorithm that follows most directly
    from re-reading the problem statement -- usually brute force or exhaustive search.
    \item \textbf{Measure why it hurts.} Compute the actual complexity, plug in the
    problem's constraints, and see the number blow up.
    \item \textbf{Find the structural insight.} Almost always this is an
    \emph{exchange argument}: show that a greedy choice can replace any part of an
    optimal solution without making it worse.
    \item \textbf{Implement the optimal algorithm}, trace it by hand, then look for the
    \emph{programming-level} tricks that make the implementation itself correct and
    fast.
  \end{enumerate}
  \begin{alertblock}{This mirrors the course}
    Steps 1--3 are exactly this week's ``exchange argument'' proof template applied to
    a new problem; step 4 is what turns a correct algorithm into working code.
  \end{alertblock}
\end{frame}

% =================================================================
\section{Problem 1 (Medium): Non-overlapping Intervals}
% =================================================================

\begin{frame}{Problem statement}
  \begin{block}{LeetCode 435 -- Non-overlapping Intervals (\textcolor{blue!50!black}{Medium})}
    Given an array of intervals \texttt{intervals[i] = [start\_i, end\_i]}, return the
    \textbf{minimum number of intervals you need to remove} so that the rest of the
    intervals are pairwise non-overlapping. Intervals that only touch at an endpoint
    (e.g.\ $[1,2]$ and $[2,3]$) do \textbf{not} count as overlapping.
  \end{block}
  \begin{itemize}
    \item Constraints (LeetCode): $1 \le n \le 10^5$.
    \item Removing the fewest intervals to leave a compatible set is the same as
    \textbf{keeping the most} intervals compatible -- i.e., this \emph{is} this week's
    Interval Scheduling problem, with the answer reported as $n$ minus the optimum.
  \end{itemize}
\end{frame}

\begin{frame}{Worked example}
  \begin{columns}
  \begin{column}{0.46\textwidth}
  \begin{center}
  \begin{tikzpicture}[>=Stealth]
    \draw[very thick] (0,3) -- (2,3) node[right,font=\scriptsize]{$[1,2]$};
    \draw[very thick] (0,2) -- (3,2) node[right,font=\scriptsize]{$[1,3]$};
    \draw[very thick, red, dashed] (2,2) -- (2,2.3);
    \draw[very thick] (2,1) -- (3,1) node[right,font=\scriptsize]{$[2,3]$};
    \draw[very thick] (3,0) -- (4,0) node[right,font=\scriptsize]{$[3,4]$};
  \end{tikzpicture}
  \end{center}
  \end{column}
  \begin{column}{0.54\textwidth}
  \small
  \texttt{intervals = [[1,2],[1,3],[2,3],[3,4]]}
  \begin{itemize}
    \item $[1,2]$, $[2,3]$, $[3,4]$ are mutually compatible (touching endpoints are
    fine) -- a compatible set of size $3$.
    \item $[1,3]$ overlaps both $[1,2]$ and $[2,3]$, so it cannot join that set.
    \item \textbf{Answer:} remove just $[1,3]$ $\Rightarrow \boxed{1}$.
  \end{itemize}
  \end{column}
  \end{columns}
\end{frame}

\begin{frame}{First instinct: search over subsets}
  \begin{block}{Most direct reading of the problem}
    ``Find the fewest removals'' $\Rightarrow$ try every subset of intervals, check
    which subsets are pairwise compatible, keep the largest compatible subset, and
    report $n$ minus its size.
  \end{block}
  \begin{algorithmic}[1]
    \State $best \gets 0$
    \For{each subset $S \subseteq \{1,\dots,n\}$}
      \If{all pairs in $S$ are pairwise compatible}
        \State $best \gets \max(best, |S|)$
      \EndIf
    \EndFor
    \State \Return $n - best$
  \end{algorithmic}
\end{frame}

\begin{frame}{Why this is bad}
  \begin{itemize}
    \item There are $2^n$ subsets, and checking pairwise compatibility of a subset of
    size $k$ costs $O(k^2)$ $\Rightarrow$ total time is exponential, roughly
    $O(2^n n^2)$ in the worst case.
    \item At $n = 10^5$, $2^n$ is astronomically larger than the number of atoms in the
    observable universe. This does not merely time out -- it is not computable in the
    lifetime of the sun.
  \end{itemize}
  \begin{alertblock}{The smell to notice}
    ``Try every subset'' is almost never the intended solution to a LeetCode Medium.
    When the problem is secretly asking for a \emph{maximum compatible subset of
    intervals}, that is a named problem this course already solved -- go find the
    lecture's algorithm instead of reinventing search.
  \end{alertblock}
\end{frame}

\begin{frame}{The structural insight}
  \begin{itemize}
    \item Minimizing removals $\equiv$ maximizing the compatible subset kept, and
    \textbf{maximizing a compatible subset of intervals is exactly this week's
    Interval Scheduling problem.}
    \item This week's exchange argument already proved: sorting by \textbf{finish
    time} and greedily keeping every interval compatible with everything kept so far
    produces a maximum compatible set.
    \item So the answer is simply $n - |\text{greedy's kept set}|$ -- no new proof
    needed, just a direct application of today's theorem.
  \end{itemize}
\end{frame}

\begin{frame}{Optimal algorithm: pseudocode}
  \begin{algorithmic}[1]
    \Function{MinRemovals}{intervals}
      \State sort intervals by end time ascending
      \State $prevEnd \gets -\infty$;\ $kept \gets 0$
      \For{$(s, e)$ in sorted intervals}
        \If{$s \ge prevEnd$}
          \State $kept \gets kept + 1$;\ $prevEnd \gets e$ \Comment{compatible: keep it}
        \EndIf
      \EndFor
      \State \Return $n - kept$
    \EndFunction
  \end{algorithmic}
  Cost: $O(n \log n)$ for the sort, $O(n)$ for the scan -- this week's exact bound.
\end{frame}

\begin{frame}{Trace on the worked example}
  \small
  \texttt{intervals = [[1,2],[1,3],[2,3],[3,4]]}, sorted by end: $[1,2],\,[1,3],\,[2,3],\,[3,4]$
  (the two end-$3$ intervals may appear in either order).
  \begin{center}
  \begin{tabular}{@{}cccl@{}}
    \toprule
    interval & $s \ge prevEnd$? & action & $prevEnd$ after \\
    \midrule
    $[1,2]$ & $1 \ge -\infty$, yes & keep (1st kept) & $2$ \\
    $[1,3]$ & $1 \ge 2$? no & \textbf{skip} (removed) & $2$ \\
    $[2,3]$ & $2 \ge 2$, yes & keep (2nd kept) & $3$ \\
    $[3,4]$ & $3 \ge 3$, yes & keep (3rd kept) & $4$ \\
    \bottomrule
  \end{tabular}
  \end{center}
  \begin{exampleblock}{Answer}
    $kept = 3$, $n = 4 \Rightarrow$ removals $= 4 - 3 = 1$ -- matches the expected
    output.
  \end{exampleblock}
\end{frame}

\begin{frame}{Complexity: naive vs.\ optimal}
  \begin{center}
  \begin{tabular}{@{}lll@{}}
    \toprule
    Approach & Time & Idea \\
    \midrule
    Search over subsets & $O(2^n n^2)$ & try every compatible subset \\
    \textbf{Sort by end time, greedy scan} & $\boldsymbol{O(n\log n)}$ & this week's
    earliest-finish-time-first \\
    \bottomrule
  \end{tabular}
  \end{center}
\end{frame}

\begin{frame}[fragile]{Python implementation}
  \begin{lstlisting}[style=py]
def eraseOverlapIntervals(intervals):
    intervals.sort(key=lambda iv: iv[1])   # sort by END time, not start time!
    prev_end = float('-inf')
    kept = 0
    for start, end in intervals:
        if start >= prev_end:              # compatible with everything kept so far
            kept += 1
            prev_end = end
    return len(intervals) - kept
  \end{lstlisting}
\end{frame}

\begin{frame}{Implementation tips \& tricks (the programming side)}
  \small
  \begin{itemize}
    \item \textbf{Sort by end time, not start time.} The single most common bug on
    this problem: sorting by \texttt{iv[0]} (start) instead of \texttt{iv[1]} (end).
    The exchange-argument proof only holds for the finish-time ordering -- swap the
    key and the algorithm silently returns a wrong (usually too-small) kept count on
    inputs where it matters.
    \item \textbf{The comparison must be \texttt{>=}, not \texttt{>}.} Touching
    endpoints ($[1,2]$ then $[2,3]$) are compatible by this problem's definition.
    Using strict \texttt{>} would incorrectly reject touching intervals -- an
    off-by-one that only shows up on boundary test cases.
    \item \textbf{Track only the last kept end time, not the whole kept set.} Because
    intervals are processed in sorted order, one running scalar (\texttt{prev\_end})
    is enough -- no need to store or re-scan the kept intervals, which is what keeps
    this $O(n)$ after the sort instead of $O(n^2)$.
  \end{itemize}
\end{frame}

% =================================================================
\section{Problem 2 (Hard): Course Schedule III}
% =================================================================

\begin{frame}{Problem statement}
  \begin{block}{LeetCode 630 -- Course Schedule III (\textcolor{red!60!black}{Hard})}
    You are given \texttt{courses[i] = [duration\_i, lastDay\_i]}: course $i$ takes
    \texttt{duration\_i} consecutive days and must be \textbf{finished} on or before day
    \texttt{lastDay\_i}. You take courses one at a time (no overlap), starting on day
    $1$. Return the \textbf{maximum number of courses} you can take.
  \end{block}
  \begin{itemize}
    \item Constraints (LeetCode): $1 \le n \le 10^4$.
    \item This is exactly this week's \textbf{job sequencing with deadlines} case
    study, generalized so each job's duration need not be $1$ -- and now the question
    is ``how many,'' not ``how much profit.''
  \end{itemize}
\end{frame}

\begin{frame}{Worked example}
  \small
  \texttt{courses = [[100,200],[1000,1250],[200,1300],[2000,3200]]}
  \begin{itemize}
    \item Sorted by deadline: $(100,200),\,(1000,1250),\,(200,1300),\,(2000,3200)$.
    \item Taking the first three in that order costs $100+1000+200 = 1300$ days total,
    which fits before the tightest of their three deadlines ($1300$).
    \item The fourth course alone needs $2000$ more days on top of that -- pushing the
    total to $3300 > 3200$, its own deadline. It cannot be added without dropping
    something.
    \item \textbf{Answer:} $3$ courses (drop the $2000$-day course).
  \end{itemize}
\end{frame}

\begin{frame}{First instinct: try every ordering}
  \begin{block}{Most direct reading of the problem}
    ``Which subset, taken in which order, fits?'' $\Rightarrow$ try every permutation
    of every subset of courses, check feasibility (running total of durations never
    exceeds the deadline of the course just placed), keep the largest feasible subset.
  \end{block}
  \begin{algorithmic}[1]
    \State $best \gets 0$
    \For{each subset $S$, each ordering $\pi$ of $S$}
      \If{$\pi$ is feasible (cumulative duration $\le$ each course's deadline in
      order)}
        \State $best \gets \max(best, |S|)$
      \EndIf
    \EndFor
    \State \Return $best$
  \end{algorithmic}
\end{frame}

\begin{frame}{Why this is bad}
  \begin{itemize}
    \item There are up to $2^n$ subsets, each with up to $n!$ orderings to check.
    \item At $n = 10^4$, this is not merely exponential -- it is
    \emph{super-exponential}, hopelessly beyond any brute-force search.
  \end{itemize}
  \begin{alertblock}{The smell to notice (again)}
    ``Try every ordering'' ignores a fact this week's lecture already proved for a
    close cousin of this problem: when jobs have deadlines, a fixed, simple ordering
    rule (sort by deadline) is enough -- you never need to consider arbitrary
    orderings.
  \end{alertblock}
\end{frame}

\begin{frame}{The structural insight}
  \begin{itemize}
    \item \textbf{Always take courses in increasing order of deadline.} Reordering two
    selected courses so the earlier-deadline one goes first can only help feasibility
    (an exchange argument identical in spirit to this week's minimizing-lateness proof)
    -- so an optimal solution can always be rearranged into deadline order without
    losing any courses.
    \item Once the order is fixed, scan courses by deadline and greedily try to take
    each one. If taking it would blow the running total past \emph{its own} deadline
    (the tightest deadline seen so far, since we're in deadline order), the schedule is
    over-subscribed -- and the best fix is to \textbf{drop whichever selected course
    has the longest duration}: that frees the most time while only affecting a course
    whose deadline is at least as loose.
    \item This is a second exchange argument: swapping the longest-duration course out
    for a shorter one never increases total time and never decreases the count of
    courses taken.
  \end{itemize}
\end{frame}

\begin{frame}{Optimal algorithm: pseudocode}
  \footnotesize
  \begin{algorithmic}[1]
    \Function{MaxCourses}{courses}
      \State sort courses by $lastDay$ ascending
      \State $H \gets$ empty max-heap (keyed by duration);\ $total \gets 0$
      \For{$(duration, lastDay)$ in sorted courses}
        \State $total \gets total + duration$;\ push $duration$ onto $H$
        \If{$total > lastDay$}
          \State $total \gets total - \Call{PopMax}{H}$ \Comment{drop the longest
          course so far}
        \EndIf
      \EndFor
      \State \Return $|H|$
    \EndFunction
  \end{algorithmic}
  Cost: $O(n \log n)$ -- one heap push/pop per course.
\end{frame}

\begin{frame}{Trace on the worked example}
  \small
  \texttt{courses = [[100,200],[1000,1250],[200,1300],[2000,3200]]}, already sorted by
  deadline.
  \begin{center}
  \begin{tabular}{@{}cccl@{}}
    \toprule
    course $(d, \ell)$ & $total$ after push & $> \ell$? & action \\
    \midrule
    $(100, 200)$ & $100$ & no & keep -- $H=\{100\}$ \\
    $(1000, 1250)$ & $1100$ & no & keep -- $H=\{100,1000\}$ \\
    $(200, 1300)$ & $1300$ & no ($1300 \le 1300$) & keep -- $H=\{100,1000,200\}$ \\
    $(2000, 3200)$ & $3300$ & \textbf{yes} ($3300 > 3200$) & pop max ($2000$),
    $total \gets 1300$ \\
    \bottomrule
  \end{tabular}
  \end{center}
  \begin{exampleblock}{Answer}
    $|H| = 3$ at the end -- the $2000$-day course is popped right back off the heap
    it was just pushed onto, since it is (unsurprisingly) the longest course seen.
  \end{exampleblock}
\end{frame}

\begin{frame}{Complexity: naive vs.\ optimal}
  \small
  \begin{center}
  \begin{tabular}{@{}lll@{}}
    \toprule
    Approach & Time & Idea \\
    \midrule
    Subsets $\times$ orderings & super-exponential & $2^n \cdot n!$ candidates \\
    \textbf{Sort by deadline + max-heap} & $\boldsymbol{O(n \log n)}$ & two exchange
    arguments \\
    \bottomrule
  \end{tabular}
  \end{center}
\end{frame}

\begin{frame}[fragile]{Python implementation}
  \begin{lstlisting}[style=py]
import heapq

def scheduleCourse(courses):
    courses.sort(key=lambda c: c[1])   # sort by deadline (lastDay)
    max_heap = []                       # Python heapq is a MIN-heap...
    total = 0
    for duration, last_day in courses:
        total += duration
        heapq.heappush(max_heap, -duration)   # ...so push -duration to simulate max
        if total > last_day:
            total += heapq.heappop(max_heap)  # pop the (negated) longest duration
    return len(max_heap)
  \end{lstlisting}
\end{frame}

\begin{frame}{Implementation tips \& tricks (the programming side)}
  \small
  \begin{itemize}
    \item \textbf{Python's \texttt{heapq} is a min-heap only.} The algorithm needs a
    \emph{max}-heap of durations. The standard trick is to push the
    \textbf{negated} duration (\texttt{-duration}); the smallest negated value is the
    largest true value. Popping then gives \texttt{-heapq.heappop(...)} back, or --
    as above -- simply \texttt{total += heapq.heappop(max\_heap)} directly, since
    adding a negative number is the same as subtracting its absolute value.
    \item \textbf{Sort key must be \texttt{lastDay} (index 1), not duration.} Sorting
    by duration instead silently breaks the exchange-argument guarantee -- the
    algorithm will still run and return \emph{a} number, just not always the right
    one, which makes this bug easy to miss on small test cases.
    \item \textbf{\texttt{len(max\_heap)} is the answer, not a separate counter.}
    Because every kept course corresponds to exactly one surviving heap entry, the
    heap's size after the scan already \emph{is} the count -- no need to maintain a
    redundant counter that could drift out of sync with the heap.
  \end{itemize}
\end{frame}

% =================================================================
\section{Summary}
% =================================================================

\begin{frame}{Summary: two problems, one lesson}
  \footnotesize
  \begin{enumerate}
    \item \textbf{Non-overlapping Intervals} is not a new problem in disguise -- it
    \emph{is} this week's Interval Scheduling problem, with the answer read off as
    $n$ minus the greedy's kept count. Recognizing the disguise is the whole
    difficulty.
    \item \textbf{Course Schedule III} extends this week's job-sequencing-with-deadlines
    case study: sort by deadline (exchange argument \#1), then use a second exchange
    argument (always evict the longest-duration course when over budget) to keep the
    schedule optimal with a single max-heap.
    \item Both times, the naive instinct was combinatorial search (subsets, or
    subsets-times-orderings); both times, an exchange argument identical in spirit to
    today's lecture collapsed the search down to one sort and one linear/heap-based
    scan.
    \item The implementation tricks (sort key choice, \texttt{>=} vs.\ \texttt{>},
    simulating a max-heap with negation) are a second, independent layer: getting the
    proof right is necessary but not sufficient -- the code has to match the proof
    exactly.
  \end{enumerate}
  \begin{alertblock}{Keep practicing}
    \texttt{LEETCODE.md} lists eight more Week 4 problems (\emph{Merge Intervals},
    \emph{Meeting Rooms II}, \emph{Gas Station}, \emph{Task Scheduler}, \dots) -- try
    applying the same four-step process to each of them.
  \end{alertblock}
\end{frame}

\end{document}

